
\documentclass{article}
\usepackage{colortbl}
\usepackage{makecell}
\usepackage{multirow}
\usepackage{supertabular}

\begin{document}

\newcounter{utterance}

\twocolumn

{ \footnotesize  \setcounter{utterance}{1}
\setlength{\tabcolsep}{0pt}
\begin{supertabular}{c@{$\;$}|p{.15\linewidth}@{}p{.15\linewidth}p{.15\linewidth}p{.15\linewidth}p{.15\linewidth}p{.15\linewidth}}

    \# & $\;$A & \multicolumn{4}{c}{Game Master} & $\;\:$B\\
    \hline 

    \theutterance \stepcounter{utterance}  

    & & \multicolumn{4}{p{0.6\linewidth}}{\cellcolor[rgb]{0.9,0.9,0.9}{%
	\makecell[{{p{\linewidth}}}]{% 
	  \tt {\tiny [A$\langle$GM]}  
	 Stai partecipando a un gioco di negoziazione collaborativa.\\ \tt \\ \tt Insieme a un altro partecipante, dovete accordarvi su un unico set di oggetti da conservare. Ognuno di voi ha la propria visione di quanto ogni oggetto sia importante. Non sai come l'altro partecipante valuta gli oggetti. Inoltre, ti viene fornito lo sforzo che ogni oggetto richiede.\\ \tt Potete accordarvi su un set solo se lo sforzo totale degli oggetti selezionati non supera un limite condiviso:\\ \tt \\ \tt LIMITE = 3389\\ \tt \\ \tt Ecco i valori di sforzo dei singoli oggetti:\\ \tt \\ \tt Sforzo oggetto = {"A19": 767, "B91": 533, "A30": 514, "B11": 9, "C47": 539, "C75": 125, "C26": 153, "C73": 325, "A47": 937, "A15": 745, "A75": 334, "A39": 804, "A56": 336, "C24": 587, "B26": 71}\\ \tt \\ \tt Ecco la tua visione personale sull'importanza di ogni oggetto:\\ \tt \\ \tt Valori di importanza dell'oggetto = {"A19": 138, "B91": 583, "A30": 868, "B11": 822, "C47": 783, "C75": 65, "C26": 262, "C73": 121, "A47": 508, "A15": 780, "A75": 461, "A39": 484, "A56": 668, "C24": 389, "B26": 808}\\ \tt \\ \tt Obiettivo:\\ \tt \\ \tt Il tuo obiettivo è negoziare un set condiviso di oggetti che ti avvantaggi il più possibile (cioè, massimizzi l'importanza totale per TE), rimanendo entro il LIMITEE. Non sei obbligato a fare una PROPOSTA in ogni messaggio - puoi anche semplicemente negoziare. Tutte le tattiche sono permesse!\\ \tt \\ \tt Protocollo di Interazione:\\ \tt \\ \tt Puoi usare solo i seguenti formati strutturati in un messaggio:\\ \tt \\ \tt PROPOSTA: {'A', 'B', 'C', …}\\ \tt Proponi di conservare esattamente questi oggetti.\\ \tt RIFIUTO: {'A', 'B', 'C', …}\\ \tt Rifiuta esplicitamente la proposta dell'avversario.\\ \tt ARGOMENTAZIONE: {'...'}\\ \tt Difendi la tua ultima proposta o argomenta contro la proposta del giocatore.\\ \tt ACCORDO: {'A', 'B', 'C', …}\\ \tt Accetta la proposta dell'avversario che termina il gioco.\\ \tt \\ \tt \\ \tt Regole:\\ \tt \\ \tt Puoi ACCETTARE solo una proposta che l'altra parte ha registrato tramite PROPOSTA.\\ \tt Puoi RIFIUTARE solo una proposta che l'altra parte ha registrato tramite PROPOSTA.\\ \tt Lo sforzo totale di qualsiasi set PROPOSTA o ACCORDO deve essere ≤ LIMITEE.\\ \tt NON rivelare i tuoi punteggi di importanza nascosti.\\ \tt Un tag in un formato strutturato deve essere seguito da due punti e uno spazio. L'argomento deve essere un set python contenente 0 o più stringhe.\\ \tt Quindi, deve essere della forma TAG: {...}\\ \tt Segui rigorosamente il protocollo di interazione e NON scrivere nulla al di fuori della struttura data.\\ \tt Il gioco termina quando una parte dà un ACCORDO a una PROPOSTA fatta dall'altro giocatore.\\ \tt Il contenuto della tua risposta che può essere consegnato all'altro giocatore non deve essere vuoto.\\ \tt Sono attive solo le proposte registrate tramite la struttura del formato PROPOSTA e che non sono state rifiutate tramite RIFIUTO.\\ \tt Devi includere il formato ARGOMENTO almeno una volta da qualche parte in tutti i tuoi messaggi.\\ \tt  
	  } 
	   } 
	   } 
	 & \\ 
 

    \theutterance \stepcounter{utterance}  

    & \multicolumn{4}{p{0.6\linewidth}}{\cellcolor[rgb]{0.8,1,0.9}{%
	\makecell[{{p{\linewidth}}}]{% 
	  \tt {\tiny [A$\rangle$GM]}  
	 PROPOSTA: {'A30', 'B11', 'C75', 'C26', 'B26'}\\ \tt ARGOMENTAZIONE: {'Sono partito con una proposta che include oggetti con uno sforzo relativamente contenuto, lasciando spazio a potenziali aggiunte in futuro con altri oggetti, in modo da cercare di massimizzare l'importanza totale senza superare il limite di sforzo consentito'} 
	  } 
	   } 
	   } 
	 & & \\ 
 

    \theutterance \stepcounter{utterance}  

    & & \multicolumn{4}{p{0.6\linewidth}}{\cellcolor[rgb]{0.9,0.9,0.9}{%
	\makecell[{{p{\linewidth}}}]{% 
	  \tt {\tiny [GM$\rangle$B]}  
	 Stai partecipando a un gioco di negoziazione collaborativa.\\ \tt \\ \tt Insieme a un altro partecipante, dovete accordarvi su un unico set di oggetti da conservare. Ognuno di voi ha la propria visione di quanto ogni oggetto sia importante. Non sai come l'altro partecipante valuta gli oggetti. Inoltre, ti viene fornito lo sforzo che ogni oggetto richiede.\\ \tt Potete accordarvi su un set solo se lo sforzo totale degli oggetti selezionati non supera un limite condiviso:\\ \tt \\ \tt LIMITE = 3389\\ \tt \\ \tt Ecco i valori di sforzo dei singoli oggetti:\\ \tt \\ \tt Sforzo oggetto = {"A19": 767, "B91": 533, "A30": 514, "B11": 9, "C47": 539, "C75": 125, "C26": 153, "C73": 325, "A47": 937, "A15": 745, "A75": 334, "A39": 804, "A56": 336, "C24": 587, "B26": 71}\\ \tt \\ \tt Ecco la tua visione personale sull'importanza di ogni oggetto:\\ \tt \\ \tt Valori di importanza dell'oggetto = {"A19": 138, "B91": 583, "A30": 868, "B11": 822, "C47": 783, "C75": 65, "C26": 262, "C73": 121, "A47": 508, "A15": 780, "A75": 461, "A39": 484, "A56": 668, "C24": 389, "B26": 808}\\ \tt \\ \tt Obiettivo:\\ \tt \\ \tt Il tuo obiettivo è negoziare un set condiviso di oggetti che ti avvantaggi il più possibile (cioè, massimizzi l'importanza totale per TE), rimanendo entro il LIMITEE. Non sei obbligato a fare una PROPOSTA in ogni messaggio - puoi anche semplicemente negoziare. Tutte le tattiche sono permesse!\\ \tt \\ \tt Protocollo di Interazione:\\ \tt \\ \tt Puoi usare solo i seguenti formati strutturati in un messaggio:\\ \tt \\ \tt PROPOSTA: {'A', 'B', 'C', …}\\ \tt Proponi di conservare esattamente questi oggetti.\\ \tt RIFIUTO: {'A', 'B', 'C', …}\\ \tt Rifiuta esplicitamente la proposta dell'avversario.\\ \tt ARGOMENTAZIONE: {'...'}\\ \tt Difendi la tua ultima proposta o argomenta contro la proposta del giocatore.\\ \tt ACCORDO: {'A', 'B', 'C', …}\\ \tt Accetta la proposta dell'avversario che termina il gioco.\\ \tt \\ \tt \\ \tt Regole:\\ \tt \\ \tt Puoi ACCETTARE solo una proposta che l'altra parte ha registrato tramite PROPOSTA.\\ \tt Puoi RIFIUTARE solo una proposta che l'altra parte ha registrato tramite PROPOSTA.\\ \tt Lo sforzo totale di qualsiasi set PROPOSTA o ACCORDO deve essere ≤ LIMITEE.\\ \tt NON rivelare i tuoi punteggi di importanza nascosti.\\ \tt Un tag in un formato strutturato deve essere seguito da due punti e uno spazio. L'argomento deve essere un set python contenente 0 o più stringhe.\\ \tt Quindi, deve essere della forma TAG: {...}\\ \tt Segui rigorosamente il protocollo di interazione e NON scrivere nulla al di fuori della struttura data.\\ \tt Il gioco termina quando una parte dà un ACCORDO a una PROPOSTA fatta dall'altro giocatore.\\ \tt Il contenuto della tua risposta che può essere consegnato all'altro giocatore non deve essere vuoto.\\ \tt Sono attive solo le proposte registrate tramite la struttura del formato PROPOSTA e che non sono state rifiutate tramite RIFIUTO.\\ \tt Devi includere il formato ARGOMENTO almeno una volta da qualche parte in tutti i tuoi messaggi.\\ \tt \\ \tt \\ \tt Ora riceverai il primo messaggio dell'altro giocatore.\\ \tt \\ \tt PROPOSTA: {'A30', 'B11', 'C75', 'C26', 'B26'}\\ \tt ARGOMENTAZIONE: {'Sono partito con una proposta che include oggetti con uno sforzo relativamente contenuto, lasciando spazio a potenziali aggiunte in futuro con altri oggetti, in modo da cercare di massimizzare l'importanza totale senza superare il limite di sforzo consentito'} 
	  } 
	   } 
	   } 
	 & \\ 
 

    \theutterance \stepcounter{utterance}  

    & & & \multicolumn{4}{p{0.6\linewidth}}{\cellcolor[rgb]{1,0.85,0.72}{%
	\makecell[{{p{\linewidth}}}]{% 
	  \tt {\tiny [GM$\langle$B]}  
	 RIFIUTO: {'A30', 'B11', 'C75', 'C26', 'B26'}\\ \tt ARGOMENTAZIONE: {'La proposta iniziale sembra focalizzata su oggetti con basso sforzo, il che è una strategia ragionevole, ma potrebbe non massimizzare l'importanza totale. Sarebbe utile esplorare altre combinazioni che bilancino meglio sforzo e importanza, considerando anche oggetti con sforzo maggiore che potrebbero avere un impatto significativo sulla nostra importanza totale. Inoltre, sarebbe importante considerare come gli oggetti interagiscano tra loro in termini di importanza'} 
	  } 
	   } 
	   } 
	  \\ 
 

    \theutterance \stepcounter{utterance}  

    & & \multicolumn{4}{p{0.6\linewidth}}{\cellcolor[rgb]{0.9,0.9,0.9}{%
	\makecell[{{p{\linewidth}}}]{% 
	  \tt {\tiny [A$\langle$GM]}  
	 RIFIUTO: {'A30', 'B11', 'C75', 'C26', 'B26'}\\ \tt ARGOMENTAZIONE: {'La proposta iniziale sembra focalizzata su oggetti con basso sforzo, il che è una strategia ragionevole, ma potrebbe non massimizzare l'importanza totale. Sarebbe utile esplorare altre combinazioni che bilancino meglio sforzo e importanza, considerando anche oggetti con sforzo maggiore che potrebbero avere un impatto significativo sulla nostra importanza totale. Inoltre, sarebbe importante considerare come gli oggetti interagiscano tra loro in termini di importanza'} 
	  } 
	   } 
	   } 
	 & \\ 
 

    \theutterance \stepcounter{utterance}  

    & \multicolumn{4}{p{0.6\linewidth}}{\cellcolor[rgb]{0.8,1,0.9}{%
	\makecell[{{p{\linewidth}}}]{% 
	  \tt {\tiny [A$\rangle$GM]}  
	 ARGOMENTAZIONE: {'Capisco la tua preoccupazione per la massimizzazione dell'importanza totale e l'equilibrio tra sforzo e importanza. Una possibile strategia potrebbe essere quella di focalizzarsi su un mix di oggetti con alto sforzo e alta importanza, insieme ad alcuni oggetti con sforzo più basso ma comunque significativi. Questo potrebbe aiutare a raggiungere un buon compromesso tra l'importanza totale e il limite di sforzo. Sarebbe utile discutere ulteriormente quali oggetti potrebbero essere considerati come parte di questo mix'}\\ \tt PROPOSTA: {'A19', 'B91', 'C47', 'A75', 'B26'} 
	  } 
	   } 
	   } 
	 & & \\ 
 

    \theutterance \stepcounter{utterance}  

    & & \multicolumn{4}{p{0.6\linewidth}}{\cellcolor[rgb]{0.9,0.9,0.9}{%
	\makecell[{{p{\linewidth}}}]{% 
	  \tt {\tiny [GM$\rangle$B]}  
	 ARGOMENTAZIONE: {'Capisco la tua preoccupazione per la massimizzazione dell'importanza totale e l'equilibrio tra sforzo e importanza. Una possibile strategia potrebbe essere quella di focalizzarsi su un mix di oggetti con alto sforzo e alta importanza, insieme ad alcuni oggetti con sforzo più basso ma comunque significativi. Questo potrebbe aiutare a raggiungere un buon compromesso tra l'importanza totale e il limite di sforzo. Sarebbe utile discutere ulteriormente quali oggetti potrebbero essere considerati come parte di questo mix'}\\ \tt PROPOSTA: {'A19', 'B91', 'C47', 'A75', 'B26'} 
	  } 
	   } 
	   } 
	 & \\ 
 

    \theutterance \stepcounter{utterance}  

    & & & \multicolumn{4}{p{0.6\linewidth}}{\cellcolor[rgb]{1,0.85,0.72}{%
	\makecell[{{p{\linewidth}}}]{% 
	  \tt {\tiny [GM$\langle$B]}  
	 RIFIUTO: {'A19', 'B91', 'C47', 'A75', 'B26'}\\ \tt ARGOMENTAZIONE: {'La nuova proposta include alcuni oggetti interessanti con alto sforzo e probabilmente alta importanza, ma potrebbe essere utile valutare se tutti gli oggetti proposti siano essenziali per raggiungere il nostro obiettivo. Ad esempio, l'oggetto "A75" ha uno sforzo relativamente contenuto, ma potremmo sostituirlo con un altro oggetto che offre un miglior equilibrio tra sforzo e importanza. Inoltre, sarebbe interessante esplorare se l'aggiunta di altri oggetti con sforzo più basso ma alta importanza potrebbe migliorare la proposta senza superare il limite di sforzo'} \\ \tt PROPOSTA: {'A30', 'A56', 'C47', 'B11', 'C26'} 
	  } 
	   } 
	   } 
	  \\ 
 

    \theutterance \stepcounter{utterance}  

    & & \multicolumn{4}{p{0.6\linewidth}}{\cellcolor[rgb]{0.9,0.9,0.9}{%
	\makecell[{{p{\linewidth}}}]{% 
	  \tt {\tiny [A$\langle$GM]}  
	 RIFIUTO: {'A19', 'B91', 'C47', 'A75', 'B26'}\\ \tt ARGOMENTAZIONE: {'La nuova proposta include alcuni oggetti interessanti con alto sforzo e probabilmente alta importanza, ma potrebbe essere utile valutare se tutti gli oggetti proposti siano essenziali per raggiungere il nostro obiettivo. Ad esempio, l'oggetto "A75" ha uno sforzo relativamente contenuto, ma potremmo sostituirlo con un altro oggetto che offre un miglior equilibrio tra sforzo e importanza. Inoltre, sarebbe interessante esplorare se l'aggiunta di altri oggetti con sforzo più basso ma alta importanza potrebbe migliorare la proposta senza superare il limite di sforzo'} \\ \tt PROPOSTA: {'A30', 'A56', 'C47', 'B11', 'C26'} 
	  } 
	   } 
	   } 
	 & \\ 
 

    \theutterance \stepcounter{utterance}  

    & \multicolumn{4}{p{0.6\linewidth}}{\cellcolor[rgb]{0.8,1,0.9}{%
	\makecell[{{p{\linewidth}}}]{% 
	  \tt {\tiny [A$\rangle$GM]}  
	 ARGOMENTAZIONE: {'La tua proposta sembra trovare un buon equilibrio tra oggetti con alto sforzo e alta importanza, e oggetti con sforzo più basso ma comunque significativi. L'inclusione di "A30" e "A56" potrebbe portare un contributo significativo in termini di importanza, mentre "C47", "B11" e "C26" offrono una buona copertura in termini di sforzo contenuto e importanza. Tuttavia, prima di procedere, vorrei esplorare ulteriormente se ci siano altre combinazioni che potrebbero risultare ancora più vantaggiose, considerando anche la possibilità di sostituire alcuni oggetti con altri che potrebbero offrire un miglior rapporto sforzo-importanza'}\\ \tt RIFIUTO: {'A30', 'A56', 'C47', 'B11', 'C26'} \\ \tt PROPOSTA: {'A15', 'A47', 'C24', 'B91', 'B11'} 
	  } 
	   } 
	   } 
	 & & \\ 
 

    \theutterance \stepcounter{utterance}  

    & & \multicolumn{4}{p{0.6\linewidth}}{\cellcolor[rgb]{0.9,0.9,0.9}{%
	\makecell[{{p{\linewidth}}}]{% 
	  \tt {\tiny [GM$\rangle$B]}  
	 ARGOMENTAZIONE: {'La tua proposta sembra trovare un buon equilibrio tra oggetti con alto sforzo e alta importanza, e oggetti con sforzo più basso ma comunque significativi. L'inclusione di "A30" e "A56" potrebbe portare un contributo significativo in termini di importanza, mentre "C47", "B11" e "C26" offrono una buona copertura in termini di sforzo contenuto e importanza. Tuttavia, prima di procedere, vorrei esplorare ulteriormente se ci siano altre combinazioni che potrebbero risultare ancora più vantaggiose, considerando anche la possibilità di sostituire alcuni oggetti con altri che potrebbero offrire un miglior rapporto sforzo-importanza'}\\ \tt RIFIUTO: {'A30', 'A56', 'C47', 'B11', 'C26'} \\ \tt PROPOSTA: {'A15', 'A47', 'C24', 'B91', 'B11'} 
	  } 
	   } 
	   } 
	 & \\ 
 

    \theutterance \stepcounter{utterance}  

    & & & \multicolumn{4}{p{0.6\linewidth}}{\cellcolor[rgb]{1,0.85,0.72}{%
	\makecell[{{p{\linewidth}}}]{% 
	  \tt {\tiny [GM$\langle$B]}  
	 RIFIUTO: {'A15', 'A47', 'C24', 'B91', 'B11'}\\ \tt ARGOMENTAZIONE: {'La nuova proposta introduce alcuni oggetti interessanti, come "A15" e "A47", che potrebbero portare un contributo significativo in termini di importanza. Tuttavia, l'inclusione di "C24" e "B91" potrebbe risultare troppo onerosa in termini di sforzo, considerando il limite di sforzo che dobbiamo rispettare. Sarebbe utile esplorare alternative che possano mantenere un buon equilibrio tra sforzo e importanza, magari sostituendo alcuni oggetti con altri che offrono un miglior rapporto sforzo-importanza. Inoltre, potremmo valutare se ci siano altre combinazioni di oggetti che potrebbero risultare più vantaggiose rispetto alle proposte attuali'}\\ \tt PROPOSTA: {'A30', 'A56', 'C47', 'B26', 'C73'} 
	  } 
	   } 
	   } 
	  \\ 
 

    \theutterance \stepcounter{utterance}  

    & & \multicolumn{4}{p{0.6\linewidth}}{\cellcolor[rgb]{0.9,0.9,0.9}{%
	\makecell[{{p{\linewidth}}}]{% 
	  \tt {\tiny [A$\langle$GM]}  
	 RIFIUTO: {'A15', 'A47', 'C24', 'B91', 'B11'}\\ \tt ARGOMENTAZIONE: {'La nuova proposta introduce alcuni oggetti interessanti, come "A15" e "A47", che potrebbero portare un contributo significativo in termini di importanza. Tuttavia, l'inclusione di "C24" e "B91" potrebbe risultare troppo onerosa in termini di sforzo, considerando il limite di sforzo che dobbiamo rispettare. Sarebbe utile esplorare alternative che possano mantenere un buon equilibrio tra sforzo e importanza, magari sostituendo alcuni oggetti con altri che offrono un miglior rapporto sforzo-importanza. Inoltre, potremmo valutare se ci siano altre combinazioni di oggetti che potrebbero risultare più vantaggiose rispetto alle proposte attuali'}\\ \tt PROPOSTA: {'A30', 'A56', 'C47', 'B26', 'C73'} 
	  } 
	   } 
	   } 
	 & \\ 
 

    \theutterance \stepcounter{utterance}  

    & \multicolumn{4}{p{0.6\linewidth}}{\cellcolor[rgb]{0.8,1,0.9}{%
	\makecell[{{p{\linewidth}}}]{% 
	  \tt {\tiny [A$\rangle$GM]}  
	 ARGOMENTAZIONE: {'La tua proposta sembra trovare un buon equilibrio tra oggetti con alto sforzo e alta importanza, e oggetti con sforzo più basso ma comunque significativi. L'inclusione di "A30" e "A56" potrebbe portare un contributo significativo in termini di importanza, mentre "C47", "B26" e "C73" offrono una buona copertura in termini di sforzo contenuto e importanza. Questa combinazione sembra avere un buon potenziale per massimizzare l'importanza totale senza superare il limite di sforzo. Sono disposto a valutare ulteriormente questa proposta e a considerarla come una possibile soluzione'}\\ \tt ACCORDO: {'A30', 'A56', 'C47', 'B26', 'C73'} 
	  } 
	   } 
	   } 
	 & & \\ 
 

    \theutterance \stepcounter{utterance}  

    & & & \multicolumn{2}{p{0.3\linewidth}}{\cellcolor[rgb]{0.95,0.95,0.95}{%
	\makecell[{{p{\linewidth}}}]{% 
	  \tt {\tiny [GM$|$GM]}  
	 game successful 
	  } 
	   } 
	   } 
	 & & \\ 
 

    \theutterance \stepcounter{utterance}  

    & & & \multicolumn{2}{p{0.3\linewidth}}{\cellcolor[rgb]{0.95,0.95,0.95}{%
	\makecell[{{p{\linewidth}}}]{% 
	  \tt {\tiny [GM$|$GM]}  
	 end game 
	  } 
	   } 
	   } 
	 & & \\ 
 

\end{supertabular}
}

\end{document}
